\setcounter{paragraph}{0}
\paragraph{Cực đại giao thoa}
\begin{itemize}
\item Những điểm cực đại giao thoa là những điểm \textit{dao động với biên độ cực đại}.
\item Đó là những điểm ứng với: 
\begin{align*}
\boder{d_1-d_2=k\lambda}\quad ; \quad (k=0;\ \pm 1;\ \pm 2;\ ...)
\end{align*}
\end{itemize}
\newghichu{
Quỹ tích của những điểm này là những đường hypebol có hai tiêu điểm là $S_1$ và $S_2$\\ 
\lefthand\ chúng được gọi là những \textit{vân giao thoa cực đại}. 
}
\paragraph{Cực tiểu giao thoa}
\begin{itemize}
\item Những điểm cực tiểu giao thoa là những điểm \textit{dao động với biên độ cực tiểu}.
\item Đó là những điểm ứng với: 
\begin{align*}
\boder{d_1-d_2=\left(k+\dfrac{1}{2}\right)\lambda}\quad ; \quad (k=0;\ \pm 1;\ \pm 2;\ ...)
\end{align*}
\end{itemize}
\newghichu{
Quỹ tích của những điểm này là những đường hypebol có hai tiêu điểm là $S_1$ và $S_2$\\ 
\lefthand\ chúng được gọi là những \textit{vân giao thoa cực tiểu}. 
}